\section{Limitations}
\label{sec:limitations}

\textbf{Domain and physical scope.} Space and time are normalized benchmark quantities. The core data use procedural bathymetry and synthetic sources; the ten GEBCO-derived crops are fully wet, depth-rescaled, and synthetically forced. These experiments assess morphology transfer rather than real-event, site-specific, or operational hazard skill.

\textbf{Reference physics vs. ground truth.} Models are trained against numerical reference solvers rather than physical observations. Discretization, boundary, and interpolation errors persist, and the Boussinesq reference uses a linear, elevation-only formulation without nonlinear wet--dry dynamics.

\textbf{Statistical replication and benchmarking.} The direct architecture comparison is based on selected seed-$18$ checkpoints, with paired intervals quantifying test-scenario variation rather than training-seed uncertainty. The seven-seed ensemble replicates only the baseline FNO and does not establish architecture-level significance. The study does not include a controlled speed comparison, so no operational latency claim is made.

\textbf{Generalization and uncertainty.} Stress tests assess domain shifts rather than out-of-family invariance. Deep ensembles remain under-dispersed, with scalar variance scaling yielding marginal coverage over space--time pixels rather than calibrated event-level probabilities.

\textbf{Data preservation and scope.} Raw trajectory mirrors, checksums, and scripts support reproducibility. This study strictly evaluates forward emulation; inverse source reconstruction and real-time data assimilation are excluded.
